\chapter{Предисловие}
Эта книжка --- про мозг как вычислительную машину. Не про то, как он выглядит и из каких долей состоит, а про то, как он работает: какие в нём элементы, как они соединены, как по соединениям идут сигналы, как машина считает, запоминает и учится и сколько энергии на это тратит. Иначе говоря, это книга по электротехнике и информатике, в которой объектом служит нейрон.

Главный герой книги --- \emph{нейрон}, и мы будем смотреть на него глазами инженера. Сначала --- как на электрическую схему: конденсатор, несколько батареек, резисторы, сопротивление которых зависит от напряжения. Мы честно, руками, выведем, почему на этой схеме в покое стоит минус семьдесят милливольт, запишем уравнения Ходжкина---Хаксли~\cite{HodgkinHuxley1952}, которые описывают нервный импульс, и увидим в них нелинейный элемент с порогом. Потом --- как на вычислительный элемент: что он складывает, что умножает, сколько бит передаёт одним импульсом, как меняет свои веса. И наконец --- как на деталь больших схем, которые фильтруют изображение, разлагают звук в спектр, управляют движением и хранят память.

Что нужно знать заранее? Немного: закон Ома, конденсатор и RC-цепь, производные и простейшие дифференциальные уравнения, логарифм и экспоненту, немного теории вероятностей и программирования. Курса физиологии не требуется. Всё биологическое, что понадобится, --- ионный канал, синапс, медиатор, --- мы будем вводить по дороге, ровно в тот момент, когда без этого нельзя двинуться дальше, и сразу переводить на язык схем и алгоритмов. Единицы выбраны те, что естественны для нейрона: напряжение в милливольтах, время в миллисекундах, длина в микрометрах, концентрации в миллимолях на литр. Мембранный потенциал всюду --- это потенциал внутри клетки минус потенциал снаружи, $V=\varphi_{\text{внутри}}-\varphi_{\text{снаружи}}$, так что у покоящегося нейрона $V<0$.

Об одной особенности построения стоит сказать сразу. На мозг можно смотреть двумя приборами. Первый --- микроэлектрод на одной клетке: он видит один нейрон во всех подробностях, каждый импульс с точностью до десятой доли миллисекунды. Второй --- электроды на коже головы или томограф: он видит миллионы клеток сразу, усреднённых по объёму размером с горошину и по времени от миллисекунд до секунд. Для инженера это знакомая пара: осциллограф на одном проводе и измеритель суммарного тока всей платы. Оба прибора смотрят на одну и ту же машину, но с разных сторон. Самые интересные места книги --- ровно те, где два прибора расходятся в показаниях, и нужно понять, почему.

Глава~\ref{sec:neurontypes}, ещё до всех частей, --- нейрон как блок с тысячами входов и одним выходом и сортировка нейронов по тому, какой сигнал идёт по их выходу. Дальше книга делится на пять частей. Часть~I (главы~\ref{sec:neuron}--\ref{sec:cable}) --- нейрон как электрическая схема: мембрана как RC-цепь с батарейками, импульс как нелинейный триггер, синапс как взвешенное соединение, дендрит и аксон как линии передачи. Часть~II (главы~\ref{sec:lif}--\ref{sec:energy}) --- нейрон как вычислитель: простые модели, вычисления в дендритах, нейронный код и информация, шум, правила обучения, энергия. Часть~III (главы~\ref{sec:circuits}--\ref{sec:oscillations}) --- схемы из нейронов: сети с обратными связями, зрение как свёртка, слух как спектральный анализ, движение как система управления, память как аттрактор, обучение с подкреплением, ритмы как тактовые генераторы. Часть~IV (главы~\ref{sec:electrophys}--\ref{sec:models}) --- измерение и моделирование: измерительные цепи для одной клетки, макроскопические измерения, численное моделирование и нейроморфное железо. Часть~V --- итоги: мозг и компьютер, пункт за пунктом.
